% Chapter 2 -- Rubric: PO1 (1.1), PO2 (research literature), feeds PO3 and PO11
\chapter{Preliminary Research}
\label{ch:research}

\section{Literature Review}
\label{sec:literature}
A formal literature review is not applicable to this project, which is an
applied engineering effort rather than a research contribution.

\section{Existing Solutions and Technologies}
\label{sec:related}
SSCL, like most software organisations, already relies on a set of
specialised tools that each cover one part of this problem: GitHub for
version control, Jira for issue tracking, SonarQube for static code
analysis, and GitHub Actions for CI/CD. Each is capable and widely adopted,
but each is scoped to its own concern and reports through its own dashboard;
none aggregates across the others, which is exactly the gap described in
Chapter 1.

A category of commercial ``engineering intelligence'' platforms --
LinearB, Swarmia, Faros AI, and Jellyfish among them -- does aggregate
metrics across the delivery pipeline, mostly throughput and cycle-time
figures pulled from version control and issue tracking. They integrate with
the same tools SSCL already uses and produce cross-tool dashboards. However,
they focus almost entirely on delivery velocity and code metrics: none
incorporates the team's own qualitative signal through anonymous surveys,
none records management actions taken in response to a decline or tracks
whether they worked, and all are priced per seat. Static-analysis and CI
dashboards (SonarCloud, GitHub Actions) remain confined to their own tool
with no cross-tool scoring at all.

The gap is a single, explainable health score spanning all four tool
categories, combined with anonymous team feedback and a record of
management action and its effectiveness.

\section{Market and User Research}
\label{sec:market}
The primary users are engineering leaders and project managers at software
organisations that run multiple concurrent projects across a mix of
specialised tools -- version control, issue tracking, static analysis, and
CI/CD -- for each one. SSCL, the project's industry partner, is
representative of this user base and was the direct source of evidence for
the need: its engineering leadership described project-health assessment as
a routine but repetitive task of opening several dashboards on a fixed
cadence and reconciling them by judgement, with no tool available to
aggregate them and no channel through which a team's own sense of how a
project was going ever reached management.

Because SSCL was the only source of user evidence available to the team, no
formal market sizing was attempted. The underlying situation, however, is not
specific to SSCL: any organisation running comparable multi-tool,
multi-project engineering operations faces the same fragmentation, which
suggests demand extends beyond the immediate industry partner.

\section{Feasibility Study}
\label{sec:feasibility}
\textbf{Technical.} GitHub, Jira, SonarCloud and GitHub Actions each expose
documented REST or GraphQL APIs with free-tier rate limits sufficient for a
handful of pilot projects, so all the data the scoring model needs was
obtainable without purchasing access. The team's existing experience with
Node.js, TypeScript, Supabase Postgres and Redis/BullMQ meant the pipeline
could be built on a stack already familiar to it. The main technical risk,
rate limits shared across projects on the same credential, was manageable by
spacing scheduled syncs. Free hosting tiers covered the API and frontend; the
one exception is the background worker, which must run continuously rather
than on demand and currently has no equivalent free serverless option, so it
remains unhosted.

\textbf{Economic.} The system was built entirely on free tiers -- Supabase,
Vercel, Upstash Redis, and the free quotas of the integrated tools -- so
development introduced no cost beyond the team's time. For SSCL, the
expected benefit is reducing the recurring manual effort of dashboard
reconciliation across its portfolio, which comes at no licensing cost to
adopt.

\textbf{Operational.} Adoption only requires the access tokens SSCL already
holds for GitHub, Jira and SonarCloud and a short one-time setup per
project; the underlying tools continue to be used exactly as before, so
operational disruption is minimal.

\textbf{Legal and regulatory.} Access tokens are stored encrypted, and
survey responses are collected with no stored link to a respondent, so no
personally identifiable data is introduced beyond ordinary account
information.

The project was judged feasible on the assumption that SSCL continues to
provide the access tokens and project structure it already uses; the
completed build, delivered within the team's two-semester timeline, confirms
that assessment.
